\documentclass[11pt]{article}
\usepackage[margin=1in]{geometry}
\usepackage{amsmath,amssymb,amsthm}
\usepackage[hidelinks]{hyperref}

\newtheorem{theorem}{Theorem}
\newtheorem{lemma}{Lemma}
\newtheorem{remark}{Remark}

\newcommand{\co}{\mathrm{co}}
\newcommand{\W}{\mathcal W}
\newcommand{\K}{\mathcal K}
\newcommand{\R}{\mathbb R}
\newcommand{\N}{\mathbb N}
\newcommand{\capdot}{\mathbin{\cap}}
\newcommand{\oplusone}{\oplus_1}

\title{A Nontrivial Tauberian Pair in \(L_1(0,1)\) Outside \(\mathcal K_+\)}
\author{FA/Banach automatic math run}
\date{2026-07-18}

\begin{document}
\maketitle

\section*{Source And Target}

The source is Manuel Gonzalez and Antonio Martinez-Abejon,
\emph{Tauberian pairs of closed subspaces of a Banach space},
arXiv:2603.27754.  For closed subspaces \(M,N\) of a Banach space \(Z\), the
source associates the pair \((M,N)\) with the operator
\[
  Q_NJ_M:M\longrightarrow Z/N.
\]
The pair is tauberian when \(Q_NJ_M\in\mathcal W_+\), and it belongs to
\(\mathcal K_+\) when \(Q_NJ_M\) is upper semi-Fredholm.

The final questions in the source ask, among other things, for non-trivial
examples in \(Z=L_1(0,1)\) of tauberian pairs
\[
  (M,N)\in \mathcal W_+\setminus \mathcal K_+.
\]
The source notes that \((M,M)\), with \(M\) an infinite-dimensional reflexive
subspace of \(L_1(0,1)\), is a trivial example.  The construction below gives
an explicit example in which both entries \(M\) and \(N\) are nonreflexive and
distinct.

\section*{The Construction}

Partition \((0,1)\) into three measurable sets \(A,B,C\) of positive measure.
Then
\[
  L_1(0,1)=L_1(A)\oplusone L_1(B)\oplusone L_1(C).
\]
Choose:
\begin{itemize}
\item an infinite-dimensional reflexive closed subspace
  \(E\subset L_1(A)\), for instance the closed span of the Rademacher
  functions on \(A\), which is isomorphic to \(\ell_2\) by Khintchine's
  inequality;
\item a closed subspace \(F\subset L_1(B)\) isometric to \(\ell_1\), for
  instance the closed span of normalized indicators of a countable measurable
  partition of \(B\);
\item a closed subspace \(G\subset L_1(C)\) isometric to \(\ell_1\), chosen
  similarly inside \(C\).
\end{itemize}
Set
\[
  M:=E\oplusone F,\qquad N:=E\oplusone G,
\]
as closed subspaces of \(Z:=L_1(0,1)\).

\begin{theorem}
The pair \((M,N)\) is a tauberian pair in \(Z=L_1(0,1)\), but
\((M,N)\notin\mathcal K_+\).  Moreover, both \(M\) and \(N\) are nonreflexive.
\end{theorem}

\begin{proof}
The spaces \(M\) and \(N\) are nonreflexive because they contain the closed
subspaces \(F\cong\ell_1\) and \(G\cong\ell_1\), respectively.

We first prove that \((M,N)\) is tauberian.  Let
\[
  P_A,\ P_B,\ P_C:L_1(0,1)\to L_1(0,1)
\]
be the band projections onto the three summands.  They extend to projections
on \(Z^{**}\), and hence induce projections
\[
  \widehat P_A,\ \widehat P_B,\ \widehat P_C: Z^\co:=Z^{**}/Z\to Z^\co
\]
whose sum is the identity.

We use the source characterization: a pair \((M,N)\) is tauberian if and only
if
\[
  M^\co\cap N^\co=\{0\}
\]
inside \(Z^\co\).  Let \(\xi\in M^\co\cap N^\co\).  Since \(M\subset
L_1(A)\oplusone L_1(B)\), the \(C\)-band component of every element of
\(M^\co\) is zero, and therefore
\[
  \widehat P_C\xi=0.
\]
Similarly, since \(N\subset L_1(A)\oplusone L_1(C)\), the \(B\)-band component
of every element of \(N^\co\) is zero, so
\[
  \widehat P_B\xi=0.
\]
It remains only to check the \(A\)-band component.  From
\(\xi\in M^\co\), the element \(\widehat P_A\xi\) belongs to the image of
\(E^{**}\) in \(Z^\co\).  But \(E\) is reflexive, so \(E^{**}=E\subset Z\),
and this image is zero in the quotient \(Z^{**}/Z\).  Hence
\[
  \widehat P_A\xi=0.
\]
Since the induced band projections sum to the identity on \(Z^\co\), we get
\(\xi=0\).  Thus \(M^\co\cap N^\co=\{0\}\), and the pair is tauberian.

Finally,
\[
  \ker(Q_NJ_M)=M\cap N=E,
\]
because the \(B\)- and \(C\)-supports are disjoint.  The space \(E\) is
infinite-dimensional.  Therefore \(Q_NJ_M\) is not upper semi-Fredholm, so
\((M,N)\notin\mathcal K_+\).
\end{proof}

\begin{remark}
This example is deliberately elementary.  It does not rely on the probabilistic
Johnson--Nasseri--Schechtman--Tkocz construction of a tauberian operator on
\(L_1(0,1)\) that is not upper semi-Fredholm.  The flexibility of pairs is
enough: the common reflexive part \(E\) makes the kernel infinite-dimensional,
while the disjoint nonreflexive \(\ell_1\) parts keep both entries of the pair
nonreflexive without creating an intersection in the quotient \(Z^\co\).
\end{remark}

\section*{Novelty Check}

The cheap run indexes were searched for arXiv:2603.27754, ``Tauberian pair'',
``closed subspaces'', ``\(L_1(0,1)\)'', ``upper semi-Fredholm'', and related
Tauberian/weak compactness keywords.  No prior run packet for this source or
for this explicit pair construction was found.  Nearby entries concern
Tauberian/semi-embedding questions for different papers, or weak compactness
questions for unrelated state-space and tensor-product settings.

\section*{Human Review Recommendation}

Review as a literal full answer to the source's request for non-trivial
examples in \(L_1(0,1)\).  The only substantive verification is the computation
\(M^\co\cap N^\co=\{0\}\), which follows from the three disjoint band
projections and the reflexivity of \(E\).

\end{document}
